\documentclass[11pt]{article}
\usepackage[T1]{fontenc}
\usepackage{lmodern}
\usepackage[margin=1in]{geometry}
\usepackage{amsmath,amssymb,amsthm,mathtools,booktabs,array}
\usepackage{xcolor,listings}
\usepackage[colorlinks=true,linkcolor=blue,citecolor=blue,urlcolor=blue]{hyperref}
\hypersetup{pdftitle={Noncategorifiability of the rank-six fusion ring B13},
 pdfsubject={An unrestricted categorification obstruction and exact certificate}}
\newtheorem{theorem}{Theorem}
\newtheorem{lemma}[theorem]{Lemma}
\newtheorem{proposition}[theorem]{Proposition}
\newtheorem{remark}[theorem]{Remark}
\newcommand{\Q}{\mathbb Q}
\newcommand{\Z}{\mathbb Z}
\newcommand{\C}{\mathbb C}
\newcommand{\Tr}{\operatorname{Tr}}
\newcommand{\Nm}{\operatorname{N}}
\newcommand{\Irr}{\operatorname{Irr}}
\newcommand{\FPdim}{\operatorname{FPdim}}
\newcommand{\OO}{\mathcal O}
\newcommand{\B}{\mathsf B_{13}}
\newcommand{\one}{\mathbf 1}
\numberwithin{equation}{section}
\lstset{basicstyle=\ttfamily\small,columns=fullflexible,keepspaces=true,
 breaklines=true,frame=single,rulecolor=\color{black!25},
 xleftmargin=5pt,xrightmargin=5pt,showstringspaces=false}
\title{Noncategorifiability of the rank-six fusion ring $\mathsf B_{13}$}
\author{}
\date{29 September 2026}
\begin{document}
\maketitle
\begin{abstract}
We exclude an explicit simple, self-dual fusion ring of rank six and multiplicity
thirteen from categorification over $\C$, without a pivotal hypothesis.
The dimension field of a hypothetical sphericalization is forced to be
$\Q(\sqrt5,\sqrt{149})$. Character orthogonality, a Casimir trace bound, and
formal-codegree divisibility reduce the possible global dimensions to a small
integer problem. A necessary norm condition for $d$-numbers leaves only the
rational value $3725$. The two complementary signed codegrees must then be roots
of $t^2-3721t+3725$, contradicting the $d$-number theorem. The only exhaustive
computation is a bounded integer check; a standard-library Python verifier and
all positive candidates are supplied.
\end{abstract}

\section{The ring and the result}
The notation $\B$ is a local label for the following ring, not a census identifier.
It has basis $1,a,b,c,d,e$, identity duality, and commutative multiplication.
Set $P=12(a+b+c+d+e)$. For nonunit basis elements $x,y$, the entry in the table
is $xy-P$; products involving $1$ are fixed by the unit rule.
\begin{center}
\small
\renewcommand{\arraystretch}{1.65}
\setlength{\tabcolsep}{4pt}
\begin{tabular}{c|ccccc}
 & $a$ & $b$ & $c$ & $d$ & $e$\\ \hline
$a$ & $1+a-2b-c+e$ & $-2a+b+c+d-e$ & $-a+b-c$ & $b$ & $a-b+e$\\
$b$ & & $1+a-3b+c+d+e$ & $a+b-2c-d+e$ & $a+b-c+e$ & $-a+b+c+d$\\
$c$ & & & $1-a-2b+c+d$ & $-b+c+d$ & $b$\\
$d$ & & & & $1+c+d+e$ & $b+d+e$\\
$e$ & & & & & $1+a+d+e$
\end{tabular}
\end{center}
All structure constants with three nonunit indices lie between $9$ and $13$. The unit coefficient
of $xy$ is $\delta_{x,y}$. The tensor is symmetric in all three indices, so
Frobenius reciprocity holds. Every nonunit square contains every basis element;
hence the ring is simple. Associativity also follows from the complete character
table below, or from the exact tensor check in the verifier.

\begin{theorem}\label{thm:main}
There is no fusion category over $\C$ whose Grothendieck ring, as a based ring,
is $\B$.
\end{theorem}

Put $s=\sqrt5$, $r=\sqrt{149}$, and $u_\pm=(61\pm5r)/2$. Two characters, in the
ordered basis $(1,a,b,c,d,e)$, are
\[
\chi_\pm=(1,u_\pm-1,u_\pm,u_\pm-1,u_\pm+1,u_\pm+1).
\]
For $v\in\{s,-s\}$, put $p=v-2$, $h=(v-1)/2$, and $k=(3-v)/2$. The other four
characters are
\[
\rho_v=(1,-p,p,h,-1,k),\qquad
\widetilde\rho_v=(1,h,p,-p,k,-1).
\]
Substitution verifies all character equations. The six rows are mutually
orthogonal and nonzero, hence independent; this also proves associativity by
identifying the complexified product with coordinatewise multiplication on
$\C^6$.
The Frobenius--Perron character is $\chi_+$. The formal codegrees are
\begin{equation}\label{eq:original-codegrees}
D,\ D',\ 25+10s,\ 25+10s,\ 25-10s,\ 25-10s,
\qquad D=\frac{18625+1525r}{2},\quad D'=\frac{18625-1525r}{2}.
\end{equation}
Writing $u=u_+$, we have
\begin{equation}\label{eq:unit}
u^2-61u=1,\qquad D=25ru,\qquad DD'=93125.
\end{equation}
For the fusion matrices $N_i=(N_{ij}^k)_{j,k}$, the Casimir trace is
\begin{equation}\label{eq:ordinary-trace}
T_0:=\Tr\left(\sum_{i=0}^5N_i^2\right)
=\sum_{i,j,k=0}^5(N_{ij}^k)^2=18725.
\end{equation}
All these identities are checked exactly from the full tensor in
\texttt{verify.py}.

\section{Reduction from an arbitrary categorification}
Assume, for contradiction, that $\mathcal C$ categorifies $\B$. Let
$\mathcal S$ be its finite sphericalization. It has two lifts $X_i^+,X_i^-$ of
each simple $X_i$, interchanged by a sign object $\varepsilon$, with
\[
\dim(X_i^-)=-\dim(X_i^+),\qquad \dim(\varepsilon)=-1,
\qquad \dim(\mathcal S)=2\dim(\mathcal C).
\]
No pivotal structure on $\mathcal C$ is needed for this construction
\cite[Corollary 7.6]{ENO4}.
Every simple of $\mathcal S$ is self-dual. Indeed, its dual lies over the same
self-dual simple of $\mathcal C$; if it were the opposite lift, sphericality
would equate opposite nonzero simple dimensions. It follows that
$K_0(\mathcal S)$ is commutative.

Choose unit-preserving lifts and form
\[
A=K_0(\mathcal S)/([\varepsilon]+1).
\]
This is free of rank six, with basis $x_0=1,x_1,\ldots,x_5$. Write
$x_ix_j=\sum_k t_{ij}^kx_k$. Then
\begin{equation}\label{eq:signed}
t_{ij}^0=\delta_{ij},\qquad
|t_{ij}^k|\le N_{ij}^k,\qquad
t_{ij}^k\equiv N_{ij}^k\pmod2.
\end{equation}
Its multiplication matrices $M_i$ are real symmetric and commute. Their faithful
regular representation is simultaneously diagonalizable, so $A_\Q$ is a
commutative semisimple algebra, with six real characters.

For any character $\lambda$ define its signed formal codegree by
\[
f_\lambda=\sum_{i=0}^5\lambda(x_i)^2.
\]
Distinct character rows are orthogonal. The standard coefficient-of-unit
pairing on $A$ has the chosen basis as an orthonormal basis; its invariance
under multiplication gives this orthogonality, equivalently the simultaneous
self-adjointness of the $M_i$. Consequently
\begin{equation}\label{eq:reciprocal}
\sum_\lambda\frac1{f_\lambda}=1,\qquad
\Tr\left(\sum_i M_i^2\right)=\sum_\lambda f_\lambda\le18725.
\end{equation}
For completeness, put $F=\operatorname{diag}(f_\lambda)$ and let $V$ be the
character matrix. Orthogonality says $VV^{\mathsf t}=F$, so
$V^{\mathsf t}F^{-1}V=I$. Its unit--unit entry is the first identity. Moreover $f_\lambda>1$: equality would force every
nonunit character value to vanish, contradicting the unit coefficient of
$x_i^2$. Each $f_\lambda$ is an algebraic integer, being an eigenvalue of the
integer matrix $\sum_iM_i^2$.

Set
\[
d_i=\dim(X_i^+),\quad d_0=1,\qquad
K=\Q(d_1,\ldots,d_5),\qquad
\Delta=\dim(\mathcal C)=\sum_{i=0}^5d_i^2.
\]
The dimension character factors through $A$. In particular $[K:\Q]\le6$.
For any original formal codegree $f$, the corresponding character pulled back
to $K_0(\mathcal S)$ has formal codegree $2f$. The central-induction theorem
\cite[Theorem 2.13]{O3} gives a central simple of dimension
\[
\frac{\dim(\mathcal S)}{2f}=\frac{\Delta}{f}.
\]
Forgetting it to $\mathcal S$ shows that
\begin{equation}\label{eq:field-divisibility}
f\in K,\qquad \Delta/f\in\OO_K.
\end{equation}
Thus $E:=\Q(s,r)\subseteq K$. Since $[E:\Q]=4$ and $[K:\Q]\le6$, we obtain
\begin{equation}\label{eq:field}
K=E=\Q(\sqrt5,\sqrt{149}).
\end{equation}
The four embeddings of $E$ therefore give four distinct characters of $A$,
with codegrees $\sigma(\Delta)$. In particular $\Delta$ is totally positive.

We also use the $d$-number theorem \cite[Theorem 1.2 and Corollary 1.4]{Od}:
$\Delta$ is a $d$-number, and so is every $f_\lambda$. For the latter assertion,
pull $\lambda$ back through $K_0(\mathcal S)\to A$; its codegree is
$2f_\lambda$. The Galois ratios of $f_\lambda$ are the same as those of
$2f_\lambda$, so the assertion follows. This step is essential: no positivity
assumption on the signed coefficients is being used.

\section{Orthogonality and the global-dimension trace}
We use the integral basis
\begin{equation}\label{eq:integer-basis}
\OO_E=\Z[\phi,\eta],\qquad
\phi=(1+s)/2,\quad \eta=(1+r)/2.
\end{equation}
Here is a direct maximality check. The product basis has discriminant
$5^2\,149^2$, so only $5$ and $149$ could divide its index in $\OO_E$.
At $5$, the polynomial $t^2-t-37$ of $\eta$ has simple roots $2,4$.
At $149$, the polynomial $t^2-t-1$ of $\phi$ has simple roots $41,109$.
Thus at either possible index prime the other quadratic factor splits, and
the localized product order is a product of the maximal integer rings in the
ramified quadratic extension. It is maximal at both primes. This proves
\eqref{eq:integer-basis}.

Write, for $1\le i\le5$,
\[
d_i=\frac{\alpha_i+\beta_i s+\gamma_i r+\delta_i sr}{4},
\qquad \alpha_i,\beta_i,\gamma_i,\delta_i\in\Z,
\]
and put $M=\sum_i\delta_i^2$. Apply orthogonality to the dimension character
and each of its three nontrivial Galois conjugates. Taking rational coefficients
in the resulting three equations gives
\begin{equation}\label{eq:equal-norms}
16+\sum_i\alpha_i^2
=5\sum_i\beta_i^2
=149\sum_i\gamma_i^2
=745\sum_i\delta_i^2=745M.
\end{equation}
In particular $M>0$, and
\begin{equation}\label{eq:trace-M}
\Tr_{E/\Q}(\Delta)=745M.
\end{equation}
The other two characters of $A$ have codegrees greater than one, so
\eqref{eq:reciprocal} gives
\[
745M+2<18725,\qquad 1\le M\le25.
\]
By \eqref{eq:unit} and \eqref{eq:field-divisibility},
\begin{equation}\label{eq:theta}
\theta:=\frac{\Delta}{25r}\in\OO_E.
\end{equation}
Using \eqref{eq:integer-basis}, write
\begin{equation}\label{eq:theta-coordinates}
\theta=\frac{x+ys+mr+wsr}{4},\qquad x,y,m,w\in\Z.
\end{equation}
Its integrality is equivalent to
\begin{equation}\label{eq:congruences}
y\equiv w\pmod2,\qquad m\equiv w\pmod2,\qquad
x\equiv y+m-w\pmod4.
\end{equation}
Since $\Tr_{E/\Q}(\Delta)=3725m$, equation \eqref{eq:trace-M} gives $M=5m$.
Hence
\begin{equation}\label{eq:m-range}
m\in\{1,2,3,4,5\},\qquad
\Delta=\frac{25}{4}\bigl(149m+149ws+xr+ysr\bigr).
\end{equation}
No reduction modulo $2$ of the fusion ring and no enumeration of signed
associative products is required.

\section{A finite norm obstruction}
\begin{lemma}[Norm--trace test]\label{lem:norm-trace}
Let $F/\Q$ be a Galois number field of degree $n$, and let $z\in\OO_F$ be a
nonzero $d$-number. Then
\[
\frac{\Tr_{F/\Q}(z)^n}{\Nm_{F/\Q}(z)}\in\Z.
\]
This remains true when $z$ generates a proper subfield of $F$.
\end{lemma}
\begin{proof}
Every $\sigma(z)/z$ is an algebraic unit by
\cite[Lemma 2.7]{Od}. Therefore
$\Tr_{F/\Q}(z)/z=\sum_\sigma\sigma(z)/z$ belongs to $\OO_F$.
Taking its norm proves the assertion.
\end{proof}

For the coordinates in \eqref{eq:m-range}, total positivity of $\Delta$ is
equivalent to
\begin{equation}\label{eq:positive}
5w^2<m^2,\qquad L>0,\qquad L^2>5H^2,
\end{equation}
where
\begin{equation}\label{eq:L-H}
L=149(m^2+5w^2)-x^2-5y^2,\qquad H=298mw-2xy.
\end{equation}
Indeed, for either embedding of $\Q(s)$, positivity at both signs of $r$
requires
\[
149(m\pm ws)>r\,|x\pm ys|.
\]
Squaring after checking the positive left sides gives \eqref{eq:positive}.
Adding the two absolute-value bounds also gives
\[
|w|\le2,\qquad |x|<mr\le5r<62,\qquad
|y|<m\sqrt{149/5}\le\sqrt{745}<28.
\]
Thus the following search bounds are proved:
\begin{equation}\label{eq:box}
-2\le w\le2,\qquad -61\le x\le61,\qquad -27\le y\le27.
\end{equation}
A relative norm calculation gives
\begin{equation}\label{eq:normtheta}
n_\theta:=\Nm_{E/\Q}(\theta)=\frac{L^2-5H^2}{256}\in\Z_{>0}.
\end{equation}
Since
\[
\Nm_{E/\Q}(\Delta)=25^4\,149^2 n_\theta,
\qquad \Tr_{E/\Q}(\Delta)=25\cdot149m,
\]
Lemma~\ref{lem:norm-trace} implies
\begin{equation}\label{eq:necessary-divisor}
n_\theta\mid149^2m^4.
\end{equation}

\begin{lemma}[Exact bounded arithmetic]\label{lem:finite}
The integer quadruples $(m,w,x,y)$ satisfying
\eqref{eq:congruences}, \eqref{eq:m-range}, \eqref{eq:positive}, and
\eqref{eq:necessary-divisor} consist of the single quadruple $(4,0,0,0)$.
Consequently $\Delta=3725$.
\end{lemma}
\begin{proof}
Use the exhaustive box \eqref{eq:box}. The numbers of integral, totally
positive candidates, and the survivors of \eqref{eq:necessary-divisor}, are
\begin{center}
\begin{tabular}{c|rrrrr|r}
$m$ & 1 & 2 & 3 & 4 & 5 & total\\ \hline
integral and totally positive & 0 & 66 & 132 & 267 & 664 & 1129\\
norm--trace survivors & 0 & 0 & 0 & 1 & 0 & 1
\end{tabular}
\end{center}
The following complete integer computation proves these counts and produces
the unique surviving quadruple. No numerical approximation is used.
\begin{lstlisting}[language=Python]
from itertools import product
counts = [0] * 6
survivors = []
for m in range(1, 6):
    for w, x, y in product(range(-2, 3), range(-61, 62),
                           range(-27, 28)):
        if 5*w*w >= m*m:
            continue
        if (y-w) % 2 or (m-w) % 2 or (x-y-m+w) % 4:
            continue
        L = 149*(m*m + 5*w*w) - x*x - 5*y*y
        H = 298*m*w - 2*x*y
        if L <= 0 or L*L <= 5*H*H:
            continue
        counts[m] += 1
        numerator = L*L - 5*H*H
        if numerator % 256:
            raise ValueError("Nonintegral norm")
        n = numerator // 256
        if (149**2 * m**4) % n == 0:
            survivors.append((m, w, x, y, n))
print(counts[1:])
print(survivors)
\end{lstlisting}
The output is
\begin{lstlisting}
[0, 66, 132, 267, 664]
[(4, 0, 0, 0, 22201)]
\end{lstlisting}
For this survivor, $\theta=r$ and $\Delta=25r^2=3725$.
The full verifier independently checks \eqref{eq:normtheta} by a
four-by-four integer determinant for each of the $1129$ candidates.
\end{proof}

\section{The two remaining codegrees give a contradiction}
We isolate the short general observation used in the final step.
\begin{lemma}[Two-codegree obstruction]\label{lem:two}
Let $f,g>0$ be algebraic integers whose unordered pair is Galois invariant,
and suppose $1/f+1/g=a/b$ in lowest terms, with $a,b$ positive integers.
Then
\[
f+g=ka,\qquad fg=kb
\]
for a positive integer $k$. If the polynomial
$t^2-kat+kb$ is irreducible over $\Q$ and $f$ is a $d$-number, then $b\mid k$.
\end{lemma}
\begin{proof}
The sum $S=f+g$ and product $P=fg$ are positive rational integers, and
$S/P=a/b$, so $(S,P)=(ka,kb)$ with $k\in\Z_{>0}$.
In the irreducible case, $f,g$ are quadratic conjugates. The $d$-number
condition implies $P\mid S^2$, by Lemma~\ref{lem:norm-trace} in that quadratic
field. Thus $kb\mid k^2a^2$, and coprimality gives $b\mid k$.
\end{proof}

\begin{proof}[Proof of Theorem~\ref{thm:main}]
By Lemma~\ref{lem:finite}, $\Delta=3725$. The four dimension characters of $A$
therefore all have codegree $3725$. Denote the other two by $f,g$.
The characteristic polynomial of the integer Casimir matrix is
\[
(t-3725)^4(t^2-St+P),\qquad S,P\in\Z,
\]
so $\{f,g\}$ is Galois invariant. By \eqref{eq:reciprocal},
\[
\frac1f+\frac1g=1-\frac4{3725}=\frac{3721}{3725}.
\]
Since $\gcd(3721,3725)=1$, Lemma~\ref{lem:two} gives
\[
S=3721k,\qquad P=3725k,\qquad k\in\Z_{>0}.
\]
The trace bound in \eqref{eq:reciprocal} gives
\[
14900+3721k\le18725,
\]
forcing $k=1$. Hence $f,g$ are the roots of
\begin{equation}\label{eq:forbidden}
t^2-3721t+3725.
\end{equation}
Its discriminant $13830941$ lies strictly between $3718^2$ and $3719^2$, so
the polynomial is irreducible. But $f$ must be a $d$-number, whereas
\[
3721^2\equiv16\pmod{3725}.
\]
This contradicts Lemma~\ref{lem:two}, equivalently the quadratic $d$-number
condition. Therefore no such $\mathcal C$ exists.
\end{proof}

\section{Certificate scope and reproduction}
The proof imports only sphericalization, the central-induction theorem, and
the $d$-number theorem. Its remaining categorical-algebra reductions are proved
above. The finite computation is not a search for associators, signed fusion
rings, or pentagon solutions. It is the bounded integer check in
Lemma~\ref{lem:finite}.

The accompanying package contains \texttt{tensor.json},
\texttt{verify.py}, \texttt{positive\_candidates.csv}, and verification reports.
The verifier uses only Python's standard library and checks the full fusion
tensor, the six characters and their codegrees, the Casimir trace, the integral
basis residue calculations, all bounded arithmetic candidates, an independent
norm determinant for each candidate, and the final quadratic contradiction.
Run
\begin{lstlisting}[language=bash]
python verify.py --output verification.json \
    --candidates positive_candidates.csv
python -O verify.py --output verification_optimized.json
\end{lstlisting}
Both commands return \texttt{PASS}. Checks use explicit exceptions and are not
disabled by \texttt{-O}. This verifies the arithmetic certificate; it is not a
formalization of the categorical inputs in a proof assistant.

\begin{thebibliography}{9}
\bibitem{ENO4}
P. Etingof, D. Nikshych, and V. Ostrik,
\emph{An analogue of Radford's $S^4$ formula for finite tensor categories},
\href{https://arxiv.org/abs/math/0404504}{arXiv:math/0404504}.
\bibitem{Od}
V. Ostrik, \emph{On formal codegrees of fusion categories},
Math. Res. Lett. \textbf{16} (2009), 895--901,
\href{https://arxiv.org/abs/0810.3242}{arXiv:0810.3242}.
\bibitem{O3}
V. Ostrik, \emph{Pivotal fusion categories of rank 3},
\href{https://arxiv.org/abs/1309.4822}{arXiv:1309.4822}.
\end{thebibliography}
\end{document}
